\section{GCN}

Let \(G=(V,E)\) have node features \(\vect{x}(v)\). Initialize
\(\vect{h}_v^{(0)}=\vect{x}(v)\), and for each iteration \(k\) update
\begin{align*}
  \vect{h}_{\mathcal N(v)}^{(k)}
    &=\operatorname{aggregate}
      \left(\left\{\vect{h}_u^{(k-1)}:u\in\mathcal N(v)\right\}\right),\\
  \vect{h}_v^{(k)}
    &=\operatorname{combine}
      \left(\vect{h}_v^{(k-1)},
            \vect{h}_{\mathcal N(v)}^{(k)}\right).
\end{align*}
The argument of \(\operatorname{aggregate}\) is a multiset. A graph-level
embedding is obtained using
\[
  \operatorname{readout}
  \left(\left\{\vect{h}_v^{(K)}:v\in V\right\}\right).
\]

\subsection{Graph Isomorphism}

Determine whether the two graphs shown in Question 1.1 of the homework PDF are
isomorphic. If they are, give a bijection between their vertices that preserves
edges. If they are not, identify a structure present in only one graph.

\begin{solution}
The two graphs are isomorphic.  Define the bijection
\[
\begin{array}{c|cccccccc}
  v & 1&2&3&4&5&6&7&8\\ \hline
  \phi(v)&A&D&E&H&C&B&G&F
\end{array}
\]
from the numbered graph to the lettered graph.  It is enough to check the
neighborhoods of nodes \(1,2,3,4\), since every edge in the numbered graph
has one endpoint in this set:
\begin{align*}
  \phi(\mathcal N(1))
    &=\phi(\{5,6,7\})=\{C,B,G\}=\mathcal N(A),\\
  \phi(\mathcal N(2))
    &=\phi(\{5,6,8\})=\{C,B,F\}=\mathcal N(D),\\
  \phi(\mathcal N(3))
    &=\phi(\{5,7,8\})=\{C,G,F\}=\mathcal N(E),\\
  \phi(\mathcal N(4))
    &=\phi(\{6,7,8\})=\{B,G,F\}=\mathcal N(H).
\end{align*}
Thus, for every pair \(u,v\),
\[
  \{u,v\}\in E_1
  \quad\Longleftrightarrow\quad
  \{\phi(u),\phi(v)\}\in E_2.
\]
Therefore \(\phi\) is an isomorphism.
\end{solution}

\subsection{Aggregation Expressiveness}

Consider the three aggregation functions
\begin{align*}
  \left[
  \operatorname{aggregate}_{\max}
  \left(\left\{\vect{h}_u^{(k-1)}:u\in\mathcal N(v)\right\}\right)
  \right]_i
  &=\max_{u\in\mathcal N(v)}
    \left[\vect{h}_u^{(k-1)}\right]_i,\\
  \operatorname{aggregate}_{\mathrm{mean}}
  \left(\left\{\vect{h}_u^{(k-1)}:u\in\mathcal N(v)\right\}\right)
  &=\frac{1}{|\mathcal N(v)|}
    \sum_{u\in\mathcal N(v)}\vect{h}_u^{(k-1)},\\
  \operatorname{aggregate}_{\mathrm{sum}}
  \left(\left\{\vect{h}_u^{(k-1)}:u\in\mathcal N(v)\right\}\right)
  &=\sum_{u\in\mathcal N(v)}\vect{h}_u^{(k-1)}.
\end{align*}

Give two graphs \(G_1,G_2\), initial node features, and nodes
\(v_1\in V_1\), \(v_2\in V_2\) such that
\(\vect{h}_{v_1}^{(0)}=\vect{h}_{v_2}^{(0)}\). After one update, the two
nodes must remain equal under mean and max aggregation but become different
under sum aggregation.

\begin{solution}
Let \(G_1\) contain one edge \(\{v_1,u\}\), and let \(G_2\) contain the two
edges \(\{v_2,a\}\) and \(\{v_2,b\}\).  Assign scalar initial features
\[
  h_{v_1}^{(0)}=h_{v_2}^{(0)}=0,
  \qquad
  h_u^{(0)}=h_a^{(0)}=h_b^{(0)}=1,
\]
and choose
\[
  \operatorname{combine}(x,y)=x+y.
\]
The neighbor-feature multisets of the distinguished nodes are
\[
  \left\{h_u^{(0)}:u\in\mathcal N(v_1)\right\}=\{1\},
  \qquad
  \left\{h_u^{(0)}:u\in\mathcal N(v_2)\right\}=\{1,1\}.
\]
Consequently,
\[
\begin{array}{c|cc}
  \text{aggregator}
    & h_{\mathcal N(v_1)}^{(1)}
    & h_{\mathcal N(v_2)}^{(1)}\\ \hline
  \max  & 1&1\\
  \mathrm{mean}&1&1\\
  \mathrm{sum}&1&2
\end{array}
\]
Since both distinguished nodes have initial feature zero, the combine step
does not change these values.  Hence max and mean give equal updated features,
whereas sum gives
\[
  h_{v_1}^{(1)}=1\ne 2=h_{v_2}^{(1)}.
\]
The example works because max and mean discard the multiplicity of repeated
neighbor features, while sum preserves it.
\end{solution}

\subsection{Relation to the Weisfeiler--Lehman Test}

At each iteration, the Weisfeiler--Lehman (WL) test updates each node
representation using its previous representation together with the multiset of
its neighbors' previous representations.

Prove that the neural graph-isomorphism test above is at most as powerful as
the WL test. More precisely, for two non-isomorphic graphs updated for \(K\)
iterations using the same \(\operatorname{aggregate}\) and
\(\operatorname{combine}\), show that
\[
  \operatorname{readout}
  \left(\left\{\vect{h}_v^{(K)}:v\in V_1\right\}\right)
  \ne
  \operatorname{readout}
  \left(\left\{\vect{h}_v^{(K)}:v\in V_2\right\}\right)
\]
implies that the WL test also declares the graphs non-isomorphic.

\begin{solution}
Let \(l_v^{(k)}\) be the WL label of node \(v\) after iteration \(k\).  The
WL hash is injective on the pair consisting of the node's previous label and
the multiset of its neighbors' previous labels.  We first prove by induction
that, for nodes from either graph,
\[
  l_u^{(k)}=l_v^{(k)}
  \quad\Longrightarrow\quad
  \vect h_u^{(k)}=\vect h_v^{(k)}.
  \tag{1}
\]

For \(k=0\), both algorithms initialize from the same node feature:
\[
  l_v^{(0)}=\vect x(v)=\vect h_v^{(0)}.
\]
Thus (1) holds at iteration zero.

Assume (1) holds after iteration \(k-1\), and suppose
\(l_u^{(k)}=l_v^{(k)}\).  Injectivity of the WL hash implies both
\[
  l_u^{(k-1)}=l_v^{(k-1)}
\]
and equality of multisets
\[
  \left\{l_a^{(k-1)}:a\in\mathcal N(u)\right\}
  =
  \left\{l_b^{(k-1)}:b\in\mathcal N(v)\right\}.
\]
By the induction hypothesis, equal previous WL labels have equal previous
neural embeddings.  Applying this implication element-by-element, with
multiplicity, gives
\[
  \left\{\vect h_a^{(k-1)}:a\in\mathcal N(u)\right\}
  =
  \left\{\vect h_b^{(k-1)}:b\in\mathcal N(v)\right\}.
\]
The shared, permutation-invariant aggregate function therefore produces the
same neighborhood embedding for \(u\) and \(v\).  Their own previous
embeddings are also equal by the induction hypothesis, so the shared combine
function gives
\[
  \vect h_u^{(k)}=\vect h_v^{(k)}.
\]
This completes the induction.

Now argue by contradiction.  Suppose the WL test does not distinguish
\(G_1\) and \(G_2\) after \(K\) iterations.  Then their final WL-label
multisets, including multiplicities, are equal:
\[
  \left\{l_v^{(K)}:v\in V_1\right\}
  =
  \left\{l_v^{(K)}:v\in V_2\right\}.
\]
By (1), a node's neural embedding is a deterministic function of its WL
label.  Hence the two final neural-embedding multisets are also equal:
\[
  \left\{\vect h_v^{(K)}:v\in V_1\right\}
  =
  \left\{\vect h_v^{(K)}:v\in V_2\right\}.
\]
Applying the same readout function to equal multisets must give equal graph
embeddings, contradicting the hypothesis that the two readouts differ.
Therefore, whenever the neural model distinguishes the graphs, the WL test
also distinguishes them.  The neural model is thus at most as powerful as
the WL test.
\end{solution}
